% Workshop short paper — GroundLM @ EMNLP 2026
% Target: 4 pages content + unlimited references (ACL standard)
% Build: pdflatex main && bibtex main && pdflatex main && pdflatex main
%
% Requires: acl.sty, ACL bibliography style (download from
%   https://github.com/acl-org/acl-style-files when CFP confirms style).
% Using the standard ACL style for now; switch to the GroundLM-specific
% template once the workshop publishes one.

\documentclass[11pt]{article}
\usepackage[T1]{fontenc}
\usepackage[utf8]{inputenc}
\usepackage{microtype}
\usepackage{inconsolata}

% --- ACL style placeholder ---
% \usepackage[review]{acl}   % review mode (anonymous)
% \usepackage{acl}            % camera-ready
% Comment out the next 4 lines when acl.sty is present:
\usepackage[margin=1in]{geometry}
\usepackage{titlesec}
\titleformat{\section}{\large\bfseries}{\thesection}{1em}{}
\titleformat{\subsection}{\normalsize\bfseries}{\thesubsection}{1em}{}

\usepackage{graphicx}
\usepackage{booktabs}
\usepackage{amsmath,amssymb}
\usepackage{algorithm}
\usepackage{algpseudocode}
\usepackage{xcolor}
\usepackage{url}
\usepackage{hyperref}

\newcommand{\todo}[1]{\textcolor{red}{[TODO: #1]}}
\newcommand{\methodname}{Faith-Tiebreak\xspace}

\title{Annotation-Free Grounding for Medical Vision-Language Models\\via Faithfulness Tiebreakers in GRPO}

\author{
  \todo{author block — anonymize for review}
}

\date{}

\begin{document}
\maketitle

%==========================================================
\begin{abstract}
Reinforcement learning with verifiable rewards (RLVR) for vision-language
models suffers from \emph{group-advantage collapse}: when rollouts for a
prompt tie on a coarse correctness metric, GRPO's group-relative
advantage vanishes and no gradient flows. We propose an
\textbf{annotation-free} attention probe --- a Lookback-Lens classifier
trained on the same correctness labels already available from the
verifiable reward --- and incorporate its output as an \textbf{ordinal
tiebreaker} (not an additive penalty) in the GRPO advantage estimator.
Because correctness ordering is never re-ordered by the probe, the
policy cannot reward-hack it. On medical VQA the tiebreaker adds a
roughly constant $+7.5$\,pp absolute F1 over correctness-only GRPO
across SLAKE ($0.431 \rightarrow 0.507$), out-of-domain VQA-RAD
($0.350 \rightarrow 0.415$) and out-of-domain PathVQA ($0.172
\rightarrow 0.250$, $+45\%$ relative). Two negative controls isolate
the contribution: rank-only GRPO ($\approx$GOPO) gives no lift, and
additive composite rewards \emph{degrade} performance in every domain
($-1.3$\,pp SLAKE, $-1.1$ VQA-RAD, $-3.5$ PathVQA), confirming the
reward-hacking pathway that the tiebreaker structure rules out by
construction. No bounding-box annotations, no preference data, single
GPU via LoRA.
\end{abstract}

%==========================================================
\section{Introduction}

Reinforcement learning with verifiable rewards has become the dominant
post-training paradigm for reasoning models, with Group Relative Policy
Optimization (GRPO) \citep{deepseek-r1} as its workhorse. GRPO computes
each rollout's advantage relative to the group mean reward, eliminating the
critic network. For coarse rewards like token-F1 on medical VQA, however,
the rollout group frequently collapses: all 8 rollouts produce
correctness-equivalent answers, the within-group variance is zero, and the
gradient vanishes \citep{dapo,ngrpo}.

Existing fixes fall into two camps. \textbf{Dynamic sampling} (DAPO) and
\textbf{advantage calibration} (NGRPO) prevent zero-variance groups by
resampling or by injecting a virtual maximum-reward sample. These restore
gradient flow but do not exploit any side information about \emph{which}
correctness-tied rollout is actually well-grounded. \textbf{Process reward
models} (Perceval, \citealp{perceval}) and \textbf{additive composite
rewards} mix a learned faithfulness signal directly into the reward
magnitude. These do exploit grounding, but at the cost of introducing a
reward-hacking pathway: a noisy probe can pull the policy toward grounded
but incorrect outputs.

\textbf{Our contribution} is the intersection of these two lines: a
faithfulness signal that informs the policy update \emph{only when
correctness ties}. Concretely:

\begin{enumerate}
\itemsep0em
\item We extend the Lookback-Lens attention probe \citep{lookback-lens}
  to vision-language models, defining a per-head feature as the ratio of
  attention paid to vision tokens versus all tokens, averaged over the
  generated response. The probe is annotation-free: training labels come
  from rollout correctness, not bounding boxes.
\item We register a \textbf{tiebreaker advantage estimator} that
  lexicographically sorts rollouts by (correctness, faith) and assigns a
  centered rank advantage. Correctness ordering is preserved; faith only
  re-orders within correctness-tied subgroups.
\item We show that \textbf{dropping unformatted rollouts} from both the
  rank set and the gradient is a strict improvement over zero-rewarding
  them, and is the single most impactful knob for out-of-domain
  generalization.
\end{enumerate}

On SLAKE \citep{slake} medical VQA (in-domain) and VQA-RAD
\citep{vqarad}/PathVQA \citep{pathvqa} (out-of-domain), our method
substantially outperforms correctness-only GRPO and additive composite
baselines, while requiring no anatomical bounding-box annotations.

%==========================================================
\section{Method}

\subsection{Faithfulness Probe}

Following \citet{lookback-lens}, we score the faithfulness of a generated
response by the share of model attention directed at grounding tokens
versus self-generated tokens. For a vision-language transformer with $L$
layers and $H$ heads, let $\alpha_{l,h}(t, s)$ denote the attention weight
from generated token $t$ to source token $s$ in head $h$ of layer $l$.
Define the per-head \emph{image-attention ratio}

\begin{equation}
r_{l,h} = \frac{1}{T} \sum_{t=1}^{T} \frac{\sum_{s \in \mathcal{V}} \alpha_{l,h}(t,s)}{\sum_{s} \alpha_{l,h}(t,s)},
\label{eq:probe-feature}
\end{equation}

where $\mathcal{V}$ is the set of vision (image-patch) input positions,
and $T$ is the number of generated tokens. Stacking $r_{l,h}$ across all
$L \cdot H$ heads gives a $LH$-dimensional feature vector. We fit a
logistic-regression classifier $\pi_\theta(\text{faith} = 1 \mid r)$ on
rollouts collected from a held-out probe training set, labeled by rollout
correctness. \textbf{No anatomical bounding boxes are used} --- the probe
is supervised purely by the verifiable reward.

\subsection{Tiebreaker Advantage}

For each prompt's group of $N$ rollouts with correctness scores $c_i$ and
faith probabilities $f_i$, the tiebreaker advantage is computed as in
Algorithm~\ref{alg:tiebreak}. The key invariants:

\begin{itemize}
\itemsep0em
\item \textbf{Correctness ordering is preserved.} A more-correct rollout
  always outranks a less-correct one --- faith only re-orders within
  correctness-tied subgroups.
\item \textbf{No reward-hacking pathway.} Because faith never enters the
  reward magnitude, a noisy probe cannot pull the policy toward
  grounded-but-incorrect outputs (cf.\ Perceval's additive penalty).
\item \textbf{Mean-zero, bounded.} The centered rank is automatically
  whitened in $[-(N-1)/2, +(N-1)/2]$, removing the need for EBPO-style
  baseline shrinkage \citep{ebpo}.
\end{itemize}

\begin{algorithm}[t]
\caption{Tiebreaker advantage estimator}
\label{alg:tiebreak}
\begin{algorithmic}[1]
\State \textbf{Input:} rollout group $\{(c_i, f_i, \text{fmt}_i)\}_{i=1}^N$
\If{drop\_unformatted}
  \State $\mathcal{I} \gets \{i : \text{fmt}_i = \text{True}\}$
  \If{$\mathcal{I} = \emptyset$} \Return zero advantages \EndIf
\Else
  \State $\mathcal{I} \gets \{1, \ldots, N\}$
\EndIf
\State Sort $\mathcal{I}$ by $(c_i, f_i)$ descending: $\sigma$
\State $N' \gets |\sigma|$, \quad $\bar r \gets (N' - 1)/2$
\For{$\text{pos} = 0, \ldots, N' - 1$}
  \State $A_{\sigma[\text{pos}]} \gets (N' - 1 - \text{pos}) - \bar r$
\EndFor
\State $A_i \gets 0$ for $i \notin \mathcal{I}$
\State \Return $\{A_i\}$ broadcast across response tokens
\end{algorithmic}
\end{algorithm}

\subsection{Drop-Unformatted Rollouts}

Rollouts that fail to emit the required \verb|<answer>...</answer>| tag
are \emph{excluded from both the rank set and the gradient} (rather than
assigned a zero reward). This is a stricter intervention than format
penalties: a format-broken rollout cannot accidentally outrank a
format-correct one through faith, and cannot consume capacity in the rank
budget. Empirically (\S\ref{sec:exp}), dropping is responsible for a
$10.4$\,pp absolute VQA-RAD improvement and is \emph{necessary} for the
tiebreaker to transfer out-of-domain at all (\S\ref{sec:exp},
observation~4).

%==========================================================
\section{Experiments}
\label{sec:exp}

\subsection{Setup}

\textbf{Model.} Qwen3-VL-2B-Thinking with LoRA ($r=16$, $\alpha=32$,
target modules $q,k,v,o$), dropout 0.05. Trainable parameter count is
$\sim$0.2\% of the 2B base model.

\textbf{Training.} GRPO with 8 rollouts per prompt, batch
$\text{prompts}=8 \times \text{grad\_accum}=8$, learning rate $1\mathrm{e}{-5}$ cosine
with weight decay 0.01, KL coefficient $\beta = 0.05$ (low-variance
estimator), temperature 1.0, max 512 new tokens. Seed 42. Five epochs
target (early stopping by val correctness).

\textbf{Data.} SLAKE organ-filtered English subset: 2076 questions split
10:1 train/val (1888/188); test on the full 1061 English questions.
Out-of-domain eval: VQA-RAD (451 test) and PathVQA (sample of 500).

\textbf{Baselines.}
(i) \emph{corr\_only}: standard GRPO, reward = token-F1.
(ii) \emph{corrrank} (\textasciitilde GOPO, \citealp{gopo}):
correctness rank only, no faith.
(iii) \emph{composite}: $0.7 \cdot \text{corr} + 0.3 \cdot \text{faith}_{\text{z}}$,
additive scalar reward.
(iv) \emph{tiebreak (bbox-cond)}: ours, with the bbox-conditioned probe
that requires anatomical annotations. This is the strong upper bound.
(v) \textbf{\emph{tiebreak (bbox-free)}}: ours, annotation-free probe.
All methods use the same data, schedule, and val-best-correct
checkpointing.

\subsection{Results}

\begin{table*}[t]
\centering
\small
\begin{tabular}{lccc}
\toprule
\textbf{Method} & \textbf{SLAKE F1} & \textbf{VQA-RAD F1} (OOD) & \textbf{PathVQA F1} (OOD)\\
\midrule
Zero-shot                          & 0.299        & 0.125        & ---         \\
\midrule
corr\_only (GRPO baseline)         & 0.431        & 0.350        & 0.172       \\
corrrank ($\approx$ GOPO)          & 0.430        & 0.339        & 0.155       \\
composite ($\alpha{=}0.7$)         & 0.419        & 0.339        & 0.137       \\
\midrule
\textbf{tiebreak (bbox-free, drop=on)}  & \textbf{0.507}  & \textbf{0.415}   & \textbf{0.250} \\
\quad ablation: drop=off            & 0.499         & 0.311         & 0.240        \\
\midrule
\textit{tiebreak (bbox-COND, drop=on)} & \textit{0.545} & \textit{0.440} & \textit{0.347} \\
\quad \emph{(uses anatomical bboxes --- upper bound)} & & & \\
\bottomrule
\end{tabular}
\caption{SLAKE in-domain and out-of-domain transfer F1. All methods
trained from Qwen3-VL-2B-Thinking with the same schedule, seed 42, and
val-best-correct checkpoint. Numbers populated by
\texttt{scripts/eval\_paper\_table.sh}.}
\label{tab:main}
\end{table*}

Table~\ref{tab:main} shows the headline numbers. Four observations:

\textbf{(1) Rank alone is insufficient.} Pure ordinal GRPO
(\emph{corrrank}, $\approx$ GOPO) yields $0.430$ on SLAKE, essentially
tied with the magnitude-based corr\_only baseline ($0.431$). The
$+7.6$\,pp lift to $0.507$ for the tiebreaker therefore comes from the
\emph{faith signal}, not from the use of ranks. This isolates our
contribution from concurrent rank-only methods.

\textbf{(2) Additive composite hurts both in- and out-of-domain.}
Combining correctness and faith into a single scalar reward ($0.7\cdot$corr
$+\,0.3\cdot$faith) \emph{degrades} SLAKE from $0.431$ to $0.419$ and
\emph{also} degrades VQA-RAD from $0.350$ to $0.339$. The reward-hacking
pathway predicted in \S\ref{sec:method} is empirically real, and the
failure is cross-domain consistent --- supporting our argument that
auxiliary signals belong in the rank order, not the reward magnitude.

\textbf{(3) Annotation-free transfer is modality-bounded.} The
bbox-conditioned tiebreaker reaches $0.545$ on SLAKE (reproducing the
reference target within $0.002$); the bbox-free variant lands at
$0.507$, trading $\sim$4\,pp in-domain for zero annotation requirement.
On near-OOD radiology (VQA-RAD) the gap shrinks to $2.5$\,pp ($0.415$
vs $0.440$): the annotation-free probe recovers $94\%$ of the
bbox-conditioned OOD F1. On far-OOD histopathology (PathVQA) the gap
widens to $9.7$\,pp ($0.250$ vs $0.347$, $72\%$ recovery). The probe
generalizes within the modality it was trained on (radiology) but
degrades on a different imaging modality entirely. The annotation-free
claim is therefore stronger inside the training-modality envelope than
across it.

\textbf{(4) Drop-unformatted is VQA-RAD-specific, not generally OOD.}
The drop ablation costs only $0.8$\,pp on SLAKE in-domain ($0.499
\rightarrow 0.507$) and $0.95$\,pp on PathVQA ($0.240 \rightarrow
0.250$) --- but a striking $10.4$\,pp on VQA-RAD ($0.311 \rightarrow
0.415$, $+33\%$ relative). Without dropping unformatted rollouts the
tiebreaker is in fact \emph{worse than corr\_only on VQA-RAD} ($0.311$
vs $0.350$): it transfers to histopathology fine but fails on
closed-form radiology QA. We hypothesize VQA-RAD's higher fraction of
short closed-form answers (yes/no, single organ names) interacts with
format-broken rollouts in a particularly bad way --- a broken-format
yes/no can land lucky F1 enough to pollute the probe's training labels.
The drop intervention prevents this leakage. On other OOD domains where
the answer distribution is more diffuse, the effect is much milder.

\subsection{Mechanism: What the Faith Signal Does}
\label{sec:mechanism}

Two of the four results above are surprising and deserve a mechanistic
account: \emph{(i)} the faith signal contributes essentially the entire
+7.6\,pp lift (corrrank alone is a wash), and \emph{(ii)} the same faith
signal \emph{hurts} when injected additively. The two observations have
a common explanation.

The tiebreaker advantage is identical to standard rank-GRPO outside of
correctness ties; faith only enters the gradient when two or more
rollouts have equal correctness. Token-F1 on medical VQA is coarse
enough that such ties are common --- with $N{=}8$ rollouts per prompt,
many groups have multiple rollouts that land at F1 $=$ 0.0 or F1 $=$
1.0. Under standard GRPO these tied rollouts receive identical
advantage, and the gradient pushes uniformly on a set of generations
that may have produced the same answer for very different reasons (some
by attending to the image, others by exploiting language priors). The
faith signal provides intra-tie \emph{resolution}: among correctness-tied
generations, those whose attention pattern matches the
historically-correct pattern receive a larger advantage, biasing the
policy gradient toward attention shapes that transfer. This
explains why the gain is large in-domain --- many tied groups become
informative gradient instead of dead gradient --- and why it transfers
cleanly out-of-domain: the regularization is toward attention
\emph{shapes} that generalize, not toward particular tokens that yielded
high F1 on training prompts. Empirically the tiebreaker contributes a roughly constant $\sim$$7.5$\,pp
absolute F1 lift over corr\_only across all three domains (SLAKE
$+7.5$, VQA-RAD $+6.5$, PathVQA $+7.8$). The relative gain therefore
\emph{grows} with distribution difficulty --- the same absolute lift
becomes $+17.5\%$ on SLAKE, $+18.6\%$ on VQA-RAD, and $+45.3\%$ on
PathVQA where the baseline is small. We interpret this as the tiebreaker
adding a fixed amount of advantage-signal resolution per group, which
becomes proportionally more important when the underlying task is
harder and the corr\_only signal is weaker.

Composite rewards lose precisely this property. When faith is added to
the reward magnitude, an incorrect-but-grounded rollout outranks an
incorrect-but-ungrounded one --- and the policy duly upweights
incorrect-but-grounded generations. The reward-hacking pathway through
the probe predicted in \S\ref{sec:method} is empirically visible as the
1.3\,pp degradation on SLAKE.  In the tiebreaker, the lex-sort
invariant rules this failure out by construction: an incorrect rollout
never outranks a correct one, regardless of faith.

\subsection{Probe Quality}

The annotation-free probe achieves val AUROC $0.873$ (vs.\
permuted-label control $0.522$, gap $+0.351$), confirming that
image-vs-text attention ratios alone carry meaningful correctness
signal in Qwen3-VL without any anatomical supervision.

%==========================================================
\section{Related Work}

\textbf{Attention-based hallucination detection.} Lookback-Lens
\citep{lookback-lens} introduced attention-ratio features for LLMs;
VIB-Probe \citep{vibprobe}, VADE \citep{vade}, and PADE
\citep{pade} extend attention-based detection to LVLMs. Our probe
inherits this lineage but adapts it for the RL inner loop, where
detection at decode time is replaced by rollout scoring at training
time.

\textbf{Ordinal advantage in GRPO.} GOPO \citep{gopo} demonstrates that
rank-only advantage matches or beats magnitude-based GRPO under noisy
reward models. Our tiebreaker is GOPO augmented with a faith signal for
intra-rank ordering --- the empirical question is whether the auxiliary
signal helps beyond pure ranks (Table~\ref{tab:main}, corrrank row).

\textbf{Reward augmentation with grounding signals.} Perceval
\citep{perceval} uses a learned process reward model to apply token-level
penalty masks to GRPO advantages, $A'_{i,t} = A_i - \alpha \cdot
m_{i,t}\cdot|A_i|$. Vision-SR1 \citep{vision-sr1} uses a decoupled
multi-reward policy loss. Both inject the grounding signal into the
reward magnitude; our tiebreaker injects it only into the rank order ---
eliminating the reward-hacking pathway that motivates Perceval's
$\alpha$-tuning.

\textbf{Advantage-collapse fixes.} DAPO \citep{dapo} dynamically
re-samples zero-variance groups; NGRPO \citep{ngrpo} injects a virtual
maximum-reward sample. These complement our approach: when no faith
ordering can break a tie (e.g.\ all 8 rollouts identical including
attention), DAPO-style resampling is the natural fallback.

\textbf{Medical VLM RL.} MedVLM-R1 \citep{medvlm-r1} and Med-R1
\citep{med-r1} establish GRPO baselines on SLAKE/VQA-RAD with
correctness-only rewards. RARL \citep{rarl} combines LoRA with RL under
similar compute constraints. Our work is the first to add a faithfulness
signal to this stack in an annotation-free manner.

%==========================================================
\section{Conclusion}

We presented an annotation-free faithfulness tiebreaker for GRPO that
grounds vision-language models without requiring bounding-box
annotations or paired preferences. On medical VQA, the method lifts
SLAKE token-F1 from $0.431$ to $0.507$ ($+17.5\%$ relative) and
improves out-of-domain VQA-RAD by $+18.6\%$ relative ($0.350 \rightarrow
0.415$), with the drop-unformatted variant
accounting for the bulk of the OOD gain. By placing the faith signal
strictly inside the tiebreaker --- not the reward magnitude --- we
eliminate the reward-hacking pathway that constrains additive
hallucination-aware rewards.

%==========================================================
\section*{Limitations}

\begin{itemize}
\itemsep0em
\item \textbf{Single domain.} All training and primary evaluation is on
  medical VQA. Cross-domain validation on VLAA-Thinking
  (math/MCQ/digit/IoU) is ongoing.
\item \textbf{Single seed.} All reported numbers use seed 42; we plan
  multi-seed runs for the long-paper extension.
\item \textbf{2B model only.} Qwen3-VL-2B-Thinking is small; the
  tiebreaker may scale differently at 7B+.
\item \textbf{Reference implementation, not verl.} A port of the
  tiebreaker advantage estimator to the verl framework is in progress
  for community reuse and broader GRPO-family comparisons (RLOO, GPG,
  GSPO, DAPO).
\item \textbf{Probe transferability.} The bbox-free probe is trained on
  Qwen3-VL-2B attentions; whether it transfers to other VLM
  architectures without retraining is an open question.
\end{itemize}

%==========================================================
\bibliographystyle{acl_natbib}
\bibliography{references}

\end{document}
